\scanpage{347}
% @ top centre, manuscript pagination
\mspage{71}

% @ upper third, crossed block, lines 1–7
\deleted{Pour terminer, je voudrais donner encore une reformulation de $A=\mathrm{Loc}(7\mathrm{a})$, très voisine de celle donnée dans le cor. 4.4.5 (page 45), où on prouve que $A$ équivaut à $C$, \uline{plus} un ``petit \unsure{chouïa|chouia}'' (pas si petit que \unsure{sa|ça} !). Ici, on va remplacer $C$ par $D$, voire par $D_{W_0}$.}

% @ heading, lines 8–9
5. \uline{Axiome des limites} $\mathrm{Loc}(8)$, \uline{et l'axiome $\mathrm{Loc}$}$(9)$ \uline{d'accessibilité}.

% @ line 10
C'est l'axiome :

% @ line 11
\uline{Loc} 8 : $W$ est stable par $\underset{\longrightarrow}{\lim}$ filtrantes.

% @ line 12
Une variante affaiblie utile est la suivante

% @ line 13
\uline{Loc} 8bis : $W$ est stable par $\underset{\longrightarrow}{\lim}$ dénombrables.

% @ lines 14–15
Explicitement, si on a \unsure{deux} systèmes inductifs filtrants dans $\deleted{\unsure{X|x}}\mathrm{Cat}$

% @ lines 16–17
\[
\begin{aligned}
\underline X\unsure{\cdot|.}&=(X_i)_{i\in I}\\
\underline Y&=(Y_i)_{i\in I},
\end{aligned}
\]

% @ line 18
et une flèche

% @ line 19
\[
\underline f=(f_i)_{i\in I}:\underline X\longrightarrow\underline Y,
\]

% @ lines 20–23
donnée par des $f_i:X_i\longrightarrow Y_i$, \deleted{alors Loc 8 nous dit que si} \inserted{posant} $X=\underset{\longrightarrow}{\lim}X_i$, $Y=\underset{\longrightarrow}{\lim}Y_i$, $f=\underset{\longrightarrow}{\lim}f_i:X\mathbin{\text{---}}Y$, Loc 8 nous dit que si les $f_i\in W$, alors $f\in W$. Ici $I$ est

% archive sheet number: 346
